\documentclass[10pt,a4paper]{amsart}
\usepackage[margin=19mm]{geometry}
\usepackage{amsmath,amssymb,mathtools}
\usepackage[scr=rsfs,cal=euler]{mathalfa}
\usepackage{xfrac}
\usepackage[unicode,hidelinks,bookmarks=false]{hyperref}
\newcommand{\ie}{i.e.\ }
\newcommand{\cf}{cf.\ }
\newcommand{\resp}{respectively\ }
\newcommand{\Subsection}{\subsection}
\providecommand{\nobreakdash}{\nobreak\hbox{-}}
\setlength{\emergencystretch}{2em}
\multlinegap0pt
\usepackage{bm}
\usepackage{bbm}
\usepackage{dsfont}
\usepackage{xspace}
\usepackage{cancel}
\usepackage{mathtools}
\usepackage{amsthm}
\makeatletter
\@ifundefined{theorem}{\newtheorem{theorem}{Theorem}}{}
\@ifundefined{proposition}{\newtheorem{proposition}[theorem]{Proposition}}{}
\makeatother
\newcommand{\Bvtex}{{B^n}}
\newcommand{\dvtex}{\,\mathrm{d}}
\newcommand{\ofpp}[1]{\operatorname{\Pi^{\ast,{\it #1}}}\!}
\newcommand{\snvtex}{{\mathbb S}^{n-1}}
\title[Affine Fractional Sobolev and Isoperimetric Inequalities]{Affine Fractional Sobolev and Isoperimetric Inequalities}
\author{Juli\'an Haddad, Monika Ludwig}
\date{Source: arXiv:2207.06375v3 (CC BY 4.0)}
\begin{document}
\maketitle

\noindent\emph{Source and licence.} Affine Fractional Sobolev and Isoperimetric Inequalities. Version arXiv:2207.06375v3 (CC BY 4.0), distributed under the Creative Commons Attribution 4.0 International licence, as recorded in the arXiv licence metadata for this exact version and shown on the abstract page https://arxiv.org/abs/2207.06375v3. Full rights notice: licenses/arxiv-2207.06375-CC-BY-4.0.txt.\par\smallskip

\noindent\textit{Editorial scope.} Counted unit: the Proposition whose source label is \texttt{prop\_Mconvexbody} in the article (source0.tex lines 794--800, source label \texttt{prop\_Mconvexbody}), delivered together with the article's own complete original proof of it (source0.tex lines 802--946, one \texttt{proof} environment of 145 lines). No mathematics is written, repaired or re-derived: statement and proof are the article's own text line for line. Required context retained uncounted: eq\_frac\_func (lines 773--789). Every definition, display and label of those lines is retained unchanged. Omitted from this package and not redistributed: 1--772, 790--793, 801--801, 947--2516. Nothing inside a retained excerpt is omitted or altered: every retained body is delivered exactly as the article's lines are written, so no omission receipt of any kind accompanies this package. Citations: the retained excerpts carry no citation command at all, so no bibliography entry of the article is transcribed and no reference is redistributed. Reference adaptation: none was needed and none was made; every reference inside the delivered unit resolves inside it, so no unresolved reference or citation is produced. Printed slips kept literal: no printed word, symbol, hypothesis or number of the counted statement or of its proof was altered. Layout and class: the article's own class is \texttt{amsart} with packages including amsmath, amsthm, amssymb, cite, mathtools, xcolor, verbatim; the wrapper is the lane's pinned \texttt{amsart} editorial wrapper at 10pt on A4, which restates the theorem-like environments the retained text needs and adds the article's own macro definitions that the retained text needs (\texttt{\textbackslash Bvtex}, \texttt{\textbackslash dvtex}, \texttt{\textbackslash ofpp}, \texttt{\textbackslash snvtex}), with extra wrapper packages (bm, bbm, dsfont, xspace, cancel, mathtools). The delivered numbering is the unit's own. Assets: no retained span contains a figure environment, a graphics-inclusion command, a PDF-page inclusion, a table or a TikZ picture; no image file is copied into this package, the package declares no asset dependency and no image file is redistributed with it. Licence: version arXiv:2207.06375v3 of this article is distributed under the Creative Commons Attribution 4.0 International licence, as recorded in the arXiv licence metadata for this exact version and shown on the abstract page https://arxiv.org/abs/2207.06375v3. A full rights notice is delivered at licenses/arxiv-2207.06375-CC-BY-4.0.txt. Characters: every retained excerpt is pure ASCII, so the delivered engine, XeLaTeX with its default Latin Modern text font, reports no missing character.
\begin{align}
\int _{\mathbb R^{n}}\int _{\mathbb R^{n}} &
\frac{|f(x)-f(y)|}{\|x-y\|_{K}^{n+s}} \dvtex x \dvtex y
\nonumber \\
&= \int _{\mathbb R^{n}}\int _{\mathbb R^{n}}
\frac{|f(y+z)-f(y)|}{\|z\|_{K}^{n+s}} \dvtex z \dvtex y
\nonumber \\
&= \int _{\snvtex} \int _{0}^{\infty }{\|t \xi \|_{K}^{-n-s}} \int _{
\mathbb R^{n}} {|f(y+t\xi )-f(y)|}\,t^{n-1} \dvtex y \dvtex t \dvtex \xi
\label{eq_frac_func}
\\
&= \int _{\snvtex}{\| \xi \|_{K}^{-n-s}} \int _{0}^{\infty }t^{-s} \Big\|
\frac{f(\cdot +t\xi ) - f(\cdot )}{t}\Big\|_{1} \dvtex t \dvtex \xi
\nonumber \\
&= \int _{\snvtex} \rho _{K}(\xi )^{n+s} \rho _{\ofpp{s} f}(\xi )^{-s}
\dvtex \xi .
\nonumber \end{align}

\begin{proposition}
\label{prop_Mconvexbody}
For non-zero $f \in W^{s,1}(\mathbb R^{n})$, the set $\,\ofpp{s} f$ is
an origin-symmetric $s$-convex body with the origin in its interior. Moreover,
there is $c>0$ depending only on $f$ such that
$\ofpp{s} f \subseteq c (1-s)^{1/s} \, \Bvtex $ for every $s \in (0,1)$.
\end{proposition}

\begin{proof}
First, note that since for $\xi \in \snvtex $ and $t>0$,
%
\begin{equation*}
\int _{\mathbb R^{n}}\vert {f(x -t\xi )-f(x)}\vert \dvtex x=\int _{
\mathbb R^{n}}\vert f(x)-f(x+t\xi )\vert \dvtex x,
\end{equation*}
%
the set $\ofpp{s} f$ is origin-symmetric.

Next, we show that $\rho _{\ofpp{s} f}$ is bounded by
$c (1-s)^{1/s} $ on $\snvtex $ for some $c>0$ depending only on $f$. We take
$r>1$ large enough so that
$\int _{r \Bvtex} \vert f(x)\vert \dvtex x \geq \frac {2}{3} \|f\|_{1}$. For given
$t \in (0,1)$, let $k>0$ be an integer such that
$2r \leq k t \leq 2r+t$. For $\xi \in \snvtex $, using changes of variables,
the triangle inequality, and the fact that $r\Bvtex $ is disjoint from
$r\Bvtex -kt\xi $, we obtain
%
%e3 #&#
%e3 #&#
%e3 #&#
%e3 #&#
%e3 #&#
%e3 #&#
%e3 #&#
\begin{align*}
\int _{\mathbb R^{n}} |f(x+t\xi ) - f(x)| \dvtex x
&= \frac {1}{k} \sum _{i=1}^{k} \int _{\mathbb R^{n}} |f(x+ i t\xi ) -
f(x+ (i-1)t \xi )| \dvtex x
\\
&\geq \frac {1}{k} \int _{\mathbb R^{n}} |f(x+ k t\xi ) - f(x)| \dvtex x \\
%\displaybreak\\
&\geq \frac {1}{k} \Big( \int _{r \Bvtex} (|f(x)| - |f(x+kt\xi )|) \dvtex x
\\
&
\hskip 52pt
+ \int _{r \Bvtex -kt\xi}( |f(x+kt\xi )| - |f(x)|) \dvtex x \Big)
\\
&\geq \frac {t}{2r+t} \Big( 2 \int _{r \Bvtex} |f(x)| \dvtex x - 2\int _{(r
\Bvtex )^{c}} |f(x)| \dvtex x \Big)
\\
&\geq \frac {2t}{2r+1} \biggl( \frac {2}{3} \|f\|_{1} - \frac {1}{3} \|f\|_{1}
\biggr).
\end{align*}
%
Hence,
%
%e3 #&#
\begin{align}
%!LEAP WARNING:  script commented \label{eq_lower_bound_norm}
\int _{0}^{\infty }t^{-s} \left \|\frac{f(\cdot +t\xi )-f(\cdot )}{t}
\right \|_{1} \dvtex t
&\geq \frac {2\|f\|_{1}}{3(2r+1)} \int _{0}^{1} t^{-s} \dvtex t
\nonumber \\
&\geq \frac {2\|f\|_{1}}{3(2r+1)} \frac {1}{1-s},
\nonumber
\end{align}
%
which implies that $\ofpp{s} f \subseteq c(1-s)^{1/s}\, \Bvtex $ for a constant
$c>0$ depending only on $f$.

Next, we show that $\Vert \cdot \Vert _{\ofpp{s}f}^{s}$ is sublinear on
$\mathbb R^{n}$. Indeed, for $\xi , \eta \in \mathbb R^{n}$, the triangle
inequality and a change of variables show that
%
%e19 #&#
\begin{equation}
%
\begin{split} \Vert \xi +\eta \Vert _{\ofpp{s} f}^{s}
&= \int \limits _{0}^{\infty }t^{-s-1} \|f(\cdot +t\xi +t\eta )-f(
\cdot )\|_{1} \dvtex t
\\
&\leq \int \limits _{0}^{\infty }t^{-s-1} ( \|f(\cdot +t\xi +t\eta )-f(
\cdot +t\xi )\|_{1} + \|f(\cdot +t\xi )-f(\cdot )\|_{1} ) \dvtex t
\label{eq_sublinearity}
\\
&= \|\xi \|_{\ofpp{s} f}^{s}+ \|\eta \|_{\ofpp{s} f}^{s},
\end{split}
%
\end{equation}
%
which shows that $\ofpp{s} f$ is an $s$-convex set.

Now, we show that $\ofpp{s} f$ has the origin in its interior. Using the
relation~\eqref{eq_frac_func} with $K=\Bvtex $, we get
%
%e3 #&#
\begin{equation}
%!LEAP WARNING:  script commented \label{eq_finite}
\int _{\snvtex} \Vert \xi \Vert _{\ofpp{s}f}^{s} \dvtex \xi = \frac {1}{n}
\int _{\mathbb R^{n}}\int _{\mathbb R^{n}}
\frac{|f(x)-f(y)|}{|x-y|^{n+s}} \dvtex x \dvtex y,
\nonumber \end{equation}
%
which is finite since $f\in W^{s,1}(\mathbb R^{n})$. We choose $r>0$ large
enough so that the set
$A = \{\xi \in \snvtex : \Vert \xi \Vert _{\ofpp{s}f}^{s} < r\}$ has positive
$(n-1)$-dimensional Hausdorff measure and a basis
$\{\xi _{1}, \ldots , \xi _{n}\} \subseteq A$ of $\mathbb R^{n}$. Writing
$x \in \mathbb R^{n}$ as $x = \sum _{i=1}^{n} x_{i} \xi _{i}$ and using~\eqref{eq_sublinearity}, we get
%
%e20 #&#
\begin{equation}
\label{eq_bounded}
%
\begin{split} \|x\|_{\ofpp sf} &\leq \Big( \sum _{i=1}^{n} |x_{i}|^{s}
\|\xi _{i}\|_{\ofpp sf}^{s} \Big)^{1/s}
\\
&\leq r^{1/s} \Big( \sum _{i=1}^{n} |x_{i}|^{s} \Big)^{1/s}
\\
&\leq r^{1/s} n^{1/s-1/2}\, \Big( \sum _{i=1}^{n} |x_{i}|^{2} \Big)^{1/2}
\\
&\leq r^{1/s} n^{1/s-1/2} \|\Xi ^{-1}\| \, |x|,
\end{split}
%
\end{equation}
%
where $\|\Xi ^{-1}\|$ is the operator norm of the inverse of the matrix
with columns $\xi _{i}$, $i=1, \ldots , n$. This shows that
$\ofpp sf$ has the origin as interior point.

To see that the gauge function is continuous combine~\eqref{eq_sublinearity} and~\eqref{eq_bounded} to obtain
%
\begin{equation*}
\|\xi +\eta \|_{\ofpp{s} f}^{s} \leq \|\xi \|_{\ofpp{s} f}^{s} + \|
\eta \|_{\ofpp{s} f}^{s} \leq \|\xi \|_{\ofpp{s} f}^{s} + d\, |\eta |^{s},
\end{equation*}
%
where $d>0$ is independent of $\xi $ and $\eta $. Applying~\eqref{eq_sublinearity} to the vectors $\xi +\eta $ and $-\eta $, we get
%
\begin{equation*}
\|\xi \|_{\ofpp{s} f}^{s} \leq \|\xi +\eta \|_{\ofpp{s} f}^{s} + \|-
\eta \|_{\ofpp{s} f}^{s},
\end{equation*}
%
which again by~\eqref{eq_bounded} implies
%
\begin{equation*}
\|\xi +\eta \|_{\ofpp{s} f}^{s} \geq \|\xi \|_{\ofpp{s} f}^{s} - d\, |
\eta |^{s}.
\end{equation*}
%
This completes the proof.
\end{proof}

\end{document}
